
\chapter{协同引导学习驱动的脑电通道选择优化方法}


\section{引言}

第三章MASER方法沿着增强的维度，从空间超分辨率重建的角度实现了从低密度到高密度脑电信号的生成式建模，有效缓解了空间分辨率受限这一技术瓶颈。作为信号级建模技术体系的重要组成部分，约简维度同样具有重要的实用意义。本章将沿着这一思路，从通道选择的角度重新审视信号级生成式建模，探讨如何在保持关键神经信息完整性的前提下，通过智能化的电极配置优化实现系统复杂度的最小化，从而构建与MASER互补协同的增强--约简技术方案。


在信号级生成建模的实践应用中，一个核心矛盾是高密度EEG系统的信息丰富性与便携式应用需求之间的冲突。随着EEG技术从实验室环境向日常应用场景的拓展\cite{casson2010wearable}，BCI系统、神经反馈训练、认知状态监测等领域对系统便携性、成本效益和实时性能提出了更高要求\cite{debener2012taking, stopczynski2014smartphone}。高密度EEG虽然能提供丰富的空间信息，但其复杂的硬件架构、高昂的成本和较差的用户体验严重限制了其在便携式场景中的部署。因此，如何在有限的硬件资源约束下最大化信号质量和任务性能，成为EEG技术产业化应用的关键挑战。

通道选择为解决这一矛盾提供了有效途径。其核心目标是从高密度电极配置中智能识别最具信息价值的电极子集，在保持神经信息完整性的同时简化系统复杂度。这一优化过程本质上是一个多目标约束的组合优化问题，需要同时考虑信号保真度、任务判别性、系统复杂度和个体适应性等相互制约的目标。现有的通道选择方法可归纳为三类主要技术路径。基于神经解剖学先验的方法通过预定义的脑区功能分工进行电极选择\cite{demuru2020comparison, luck2000event, nomenclature1991american}，如运动想象任务中的感觉运动皮层电极选择。这类方法具有明确的理论依据和良好的可解释性，但难以适应神经活动的分布式特性、个体差异和任务特异性需求。数据驱动的统计方法通过互信息、信噪比、功率谱密度等指标进行通道排序\cite{handiru2016optimized, yang2024improving}，能够自动发现数据中的统计规律，但往往基于单一评价准则，难以平衡多重优化目标。基于深度学习的方法通过可学习的选择机制联合优化任务性能和电极选择\cite{craik2019deep, zhang2021motor, alotaiby2015review}，具有强大的非线性建模能力，但容易过拟合且缺乏在全局信息保持与任务特征提取之间的理论平衡机制。


为解决上述技术挑战，本章提出了基于协同引导学习的通道选择优化方法CogSSM（\textbf{Co}-\textbf{G}uidance learning with \textbf{S}tate \textbf{S}pace \textbf{M}odeling）。该方法的设计灵感源于认知神经科学中大脑信息处理的双重机制：全局整合与任务特化\cite{bertolero2015modular, ito2020cortical}。大脑通过分布式神经网络维持对环境的整体感知和内在状态的统一表征，同时特定脑区能够针对当前任务进行专门化处理\cite{majovski1982cognitive, kanwisher2010functional}。CogSSM将这一认知原理转化为可计算的算法框架，设计了双路径协同优化架构。任务无关路径采用超分辨率重建目标\cite{zhang2024maser}，模拟大脑的分布式信息处理机制，确保在降维配置下维持全局EEG特征的统计一致性和频谱完整性；任务相关路径通过特定的分类或回归目标优化任务性能，模拟大脑的专门化处理机制，优先保留对目标任务至关重要的判别性特征。通过状态空间建模技术的引入，CogSSM能够高效处理EEG信号的长时序特性，实现计算效率与建模能力的良好平衡。

本章将详细阐述CogSSM方法的理论基础、双路径协同机制的设计原理，以及状态空间建模在通道选择任务中的具体应用，通过实验验证其在多种EEG应用场景中的有效性和泛化能力。


\section{脑电通道选择优化方法}
\subsection{问题定义与优化目标}
通道选择本质上可视为EEG空间超分辨率的逆向过程，即从高密度空间配置向低密度配置的约简。这一过程涉及对高维EEG信号中信息分布规律的深层理解，以及不同电极位置对特定认知任务贡献度的精确量化\cite{alotaiby2015review}。
设原始高密度EEG信号为$\mathbf{X}_h \in \mathbb{R}^{C_h \times T}$，其中$C_h$表示电极总数，$T$表示时间采样点数。最优通道选择的目标是从$C_h$个电极中识别出$C_l$个电极（$C_l < C_h$），使得选择后的低维信号$\mathbf{X}_l \in \mathbb{R}^{C_l \times T}$能够保留原始信号中的关键神经信息\cite{alyasseri2020person}。具体而言，降维后的信号可表示为：
\begin{equation}
    \mathbf{X}_l = \mathbf{S} \mathbf{X}_h,
\end{equation}
式中，$\mathbf{S} \in \{0, 1\}^{C_l \times C_h}$是二值选择矩阵，满足约束条件$\|\mathbf{S} \mathbf{1}_{C_h}\|_1 = C_l$以确保恰好选择$C_l$个电极，$\mathbf{1}_{C_h}$是$C_h$维的全1向量。

通道选择问题的核心在于学习最优的选择矩阵$\mathbf{S}$，使其在最大化下游任务性能的同时最小化所选通道间的冗余性。该优化问题可形式化表述为：
\begin{equation}
    \mathbf{S}^* = \arg\max_{\mathbf{S}} \mathcal{Q}\big(\mathbf{X}_l(\mathbf{S})\big) \quad \text{s. t.} \quad \|\mathbf{S} \mathbf{1}_{C_h}\|_1 = C_l,
\end{equation}
式中，$\mathcal{Q}(\cdot)$表示使用降维信号$\mathbf{X}_l(\mathbf{S})$在特定任务上的性能评估函数。这一优化目标体现了通道选择方法需要在信号保真度与任务特异性之间取得平衡，既要保持足够的神经信息用于下游分析，又要去除冗余和噪声成分。单目标优化策略往往难以充分权衡信息完整性与任务特异性之间的复杂关系。为解决这一理论与实践难题，本研究提出CogSSM方法，采用协同引导学习策略构建双重优化机制。该方法通过任务无关的全局信息保持分支与任务相关的特征优化分支，在信号重构质量和下游任务性能之间建立动态平衡，从而实现更加鲁棒且具有良好泛化能力的通道选择效果。


\subsection{CogSSM网络架构设计}
CogSSM方法的核心架构如图~\ref{fig:C4-1}所示，该框架通过三个关键组件的协同作用实现自适应通道选择：自适应通道选择模块、ConvSS2D特征提取主干网络以及协同引导学习策略。整体流程详见算法\ref{alg:cogssm}。

\begin{figure*}[htbp]
    \centering
    \includegraphics[width=\linewidth]{C4-1框架-CN.pdf}
    \caption{CogSSM通道选择模型整体框架}
    \label{fig:C4-1}
\end{figure*}

适应通道选择模块构成了整个框架的前端处理单元，其采用Gumbel-Softmax重参数化技巧解决离散选择优化的不可微分性问题。该模块首先通过重要性评分网络对输入的高密度EEG信号$\mathbf{X}_h$计算各通道的信息价值度量，生成重要性分布$\mathbf{p} \in \mathbb{R}^{C_h}$。随后，利用Gumbel-TopK采样机制从该分布中选择$C_l$个最具判别性的通道，形成可微分的二值选择矩阵$\mathbf{S}_b \in \{0,1\}^{C_l \times C_h}$。这一设计使得离散的通道选择过程能够嵌入到连续的端到端优化框架中。ConvSS2D特征提取主干网络负责从选择后的低维信号$\mathbf{X}_l$中提取高质量的神经表征。该网络采用分层特征学习架构，首先通过卷积层捕获EEG信号的局部时空模式和频域特性，然后利用SS2D块建模长程时序依赖关系。SS2D块通过多方向扫描机制同时处理时间和空间维度的信息流动，能够有效捕获EEG信号中复杂的神经动力学特征和跨通道耦合模式。这种设计充分考虑了EEG信号的生理学特性，即神经活动在时间上的连续性和空间上的分布式特征。
协同引导学习策略是CogSSM方法的核心贡献，通过双分支架构实现任务特异性与信号完整性的动态平衡。任务感知分支专注于下游应用的性能优化，将ConvSS2D提取的特征输入分类器，通过交叉熵损失$\mathcal{L}_{\text{cls}}$最大化任务相关的判别能力。任务无关分支则致力于维持全局信号结构的完整性，通过超分辨率重建模块将降维特征恢复为原始高密度信号$\hat{\mathbf{X}}_h$，并采用重构损失$\mathcal{L}_{\text{recon}}$确保信息保真度。此外，引入多样性正则化项$\mathcal{L}_{\text{div}}$防止选择策略过度集中于少数通道，促进选择结果的空间分布合理性。

\begin{algorithm}[htbp]
\caption{\textbf{CogSSM通道选择算法流程}}
\label{alg:cogssm}
\KwIn{高密度EEG信号 $\mathbf{X}_h$: {\texttt{(C\textsubscript{h}, T)}}；目标通道数 $C_l$；类别数 $N_c$}
\KwOut{最优选择矩阵 $\mathbf{S}^*$: {\texttt{(C\textsubscript{l}, C\textsubscript{h})}} }
\SetKwComment{Comment}{$\triangleright$\ }{}

\Comment{基于Gumbel-TopK的自适应通道选择}
\Begin{
    $\mathbf{p}$: {\texttt{(C\textsubscript{h})}} $\leftarrow \text{ScoreImportance}(\mathbf{X}_h)$\;
    
    $\mathbf{S}_b$: {\texttt{(C\textsubscript{l}, C\textsubscript{h})}} $\leftarrow \text{GumbelSoftmax}(\mathbf{p}, \tau)$\;
    
    $\mathbf{X}_l$: {\texttt{(C\textsubscript{l}, T)}} $\leftarrow \text{Tokenize}(\mathbf{X}_h, \mathbf{S}_b)$\;
} 

\Comment{ConvSS2D特征提取与嵌入}
\Begin{
    $\mathbf{E}_{\text{embed}}$: {\texttt{(N, E)}} $\leftarrow \text{SelectiveEmbed}(\mathbf{X}_l)$\;
    
    $\mathbf{E}$: {\texttt{(N, E)}} $\leftarrow \text{ConvSS2DEnc}(\mathbf{E}_{\text{embed}})$\;
}

\Comment{双分支协同引导学习}
\Begin{
    $\mathbf{y}$: {\texttt{(N\textsubscript{c})}} $\leftarrow \text{Classifier}(\mathbf{E})$\;
    
    $\hat{\mathbf{X}}_l$: {\texttt{(C\textsubscript{l}, T)}} $\leftarrow \text{EmbedInverse}(\mathbf{E})$\;
    
    $\hat{\mathbf{X}}_l$: {\texttt{(C\textsubscript{l}, T)}} $\leftarrow \text{AdjustTime}(\hat{\mathbf{X}}_l)$\;
    
    $\hat{\mathbf{X}}_h$: {\texttt{(C\textsubscript{h}, T)}} $\leftarrow \text{SuperResolution}(\hat{\mathbf{X}}_l, \mathbf{S}_b)$\;
}

\Comment{联合损失优化，通过$\nabla_\theta \mathcal{L}_{\text{sel}}$更新参数$\theta$}
\Begin{
    $\mathcal{L}_{\text{recon}} \leftarrow \|\mathbf{X}_h - \hat{\mathbf{X}}_h\|_2^2$\;
    
    $\mathcal{L}_{\text{cls}} \leftarrow \text{CrossEntropy}(\mathbf{y}, \mathbf{y}_{\text{true}})$\;
    
    $\mathcal{L}_{\text{div}} \leftarrow \text{DiversityLoss}(\mathbf{S}_b)$\;
    
    $\mathcal{L}_{\text{sel}} \leftarrow \mathcal{L}_{\text{recon}} + \lambda_{\text{cls}}\mathcal{L}_{\text{cls}} + \lambda_{\text{div}}\mathcal{L}_{\text{div}}$\;
}

$\mathbf{S}^* \leftarrow \arg\min_{\mathbf{S}} \{ \mathcal{L}_{\text{sel\_val}}\}$\;

\Return $\mathbf{S}^*$\;
\end{algorithm}

\subsection{可微分离散通道选择机制}
通道选择本质上是一个组合优化问题，需要从$C_h$个候选电极中选择$C_l$个最优子集。这类离散优化问题的搜索空间呈指数级增长，当电极数量为64、需要选择20个时，可能的组合数量达到$\binom{64}{20} \approx 4.19 \times 10^{15}$，使得穷举搜索在计算上完全不可行。传统的离散优化方法如贪婪搜索、遗传算法、模拟退火等虽然能够处理这类问题，但存在计算复杂度高、容易陷入局部最优等限制。离散选择操作的不可微性阻碍了梯度信息的传播，使得无法使用基于梯度的优化算法。在深度学习框架中，反向传播算法依赖于损失函数相对于网络参数的梯度信息来更新参数。然而，离散的argmax或top-k操作会导致梯度为零或未定义，从而无法进行有效的参数更新。这种不可微性问题是将通道选择集成到端到端深度学习系统中的根本障碍。
为了解决这一核心问题，CogSSM框架采用Gumbel-Softmax重参数化技术来实现可微分的离散选择\cite{strypsteen2021end}。该技术通过引入Gumbel噪声和温度参数来近似离散分布的采样过程，从而在保持选择语义的同时实现梯度的可传播性。


Gumbel-Softmax方法的理论基础来源于Gumbel分布的性质。对于离散分布${p_1, p_2, \ldots, p_n}$，其中$\sum_{i=1}^n p_i = 1$，标准的离散采样过程为：
\begin{equation}
z = \text{one-hot}(\arg\max_i \log p_i),
\end{equation}
式中，$z$表示one-hot编码的选择向量，$p_i$表示第$i$个候选项的选择概率。这个过程是完全不可微的。Gumbel-Softmax通过以下变换实现可微近似：
\begin{equation}
z_i = \frac{\exp((\log p_i + g_i)/\tau)}{\sum_{j=1}^n \exp((\log p_j + g_j)/\tau)},
\end{equation}
式中，$z_i$表示第$i$个元素的连续近似值，$g_i = -\log(-\log(u_i))$表示从标准Gumbel分布采样的噪声项，$u_i \sim \text{Uniform}(0,1)$表示均匀分布随机变量，$\tau > 0$表示温度参数。当$\tau \rightarrow 0$时，该分布趋近于one-hot分布；当$\tau \rightarrow \infty$时，分布变为均匀分布。

针对通道选择的特殊需求，本研究扩展了标准Gumbel-Softmax来实现Top-K选择而非单元素选择。自适应通道选择模块首先通过分层卷积网络进行时序特征提取，计算通道级重要性分数：
\begin{equation}
\mathbf{f} = \text{MLP}(\text{Conv1D}(\mathbf{X}_h)) \in \mathbb{R}^{B \times C_h},
\end{equation}
式中，$\mathbf{f}$表示原始通道重要性分数，$B$表示批次大小，$C_h$表示总通道数，MLP表示由两个线性变换层和LeakyReLU激活函数组成的多层感知机。


稳定化选择机制通过引入历史信息来平滑选择过程：
\begin{equation}
\hat{\mathbf{f}} = \mathbf{f} + \gamma\mathbf{h}_{t-1},
\end{equation}
式中，$\hat{\mathbf{f}}$表示稳定化后的重要性分数，$\gamma$表示历史信息的影响系数，$\mathbf{h}_{t-1}$表示前一时刻的历史记录。历史记录$\mathbf{h}_t \in \mathbb{R}^{C_h}$维护训练迭代过程中通道选择频率：
\begin{equation}
\mathbf{h}_t = \alpha\mathbf{h}_{t-1} + (1-\alpha)\mathbb{E}[\mathbf{S}_t],
\end{equation}
式中，$\alpha$表示动量系数，$\mathbb{E}[\mathbf{S}_t]$表示当前时刻选择矩阵的期望值。二值选择矩阵$\mathbf{S} \in \{0,1\}^{B \times C_l \times C_h}$通过可微分的Gumbel-TopK采样过程生成：
\begin{equation}
\mathbf{S}_{b,i,j} = \begin{cases}
1 & \text{if } j \in \text{top}_k(\mathbf{p}_b), \\
0 & \text{其他},
\end{cases}
\end{equation}
式中，$\mathbf{S}_{b,i,j}$表示批次$b$中第$i$个选择位置对第$j$个通道的选择状态，$\text{top}_k(\mathbf{p}_b)$表示概率$\mathbf{p}_b$中前$k$个最大值的索引集合。选择概率$\mathbf{p}_b$计算为：
\begin{equation}
\mathbf{p}_b = \text{softmax}\left(\frac{\hat{\mathbf{f}}_b + \mathbf{g}_b}{\tau}\right),
\end{equation}
式中，$\mathbf{g}_b = -\log(-\log(\mathbf{u}_b))$表示从$\mathbf{u}_b \sim \text{Uniform}(0,1)$抽取的标准Gumbel噪声样本，$\tau$表示温度参数。温度参数遵循指数衰减调度，逐步锐化概率分布：
\begin{equation}
\tau = \tau_{\text{init}}\exp(-t/\beta),
\end{equation}
式中，$\tau_{\text{init}}$表示初始温度值，$t$表示当前训练步数，$\beta$表示退火速率控制参数。

稳定化机制的设计旨在解决传统Gumbel-Softmax方法在EEG通道选择任务中的固有局限性。通过历史信息的有机融合，该机制有效避免了训练过程中由于随机性导致的选择不稳定问题，同时保持了对新兴模式的学习能力。温度参数的动态调整策略进一步确保了网络在探索阶段能够充分考虑不同的通道组合方案，并在收敛阶段锁定最优的选择配置，最终实现高质量且稳定的通道选择效果。

\subsection{ConvSS2D特征提取模块设计}
状态空间模型（State Space Models，SSMs）在序列建模方面表现出色，但往往忽略了多通道EEG信号中的关键空间模式。受到VMamba中VSS块的启发\cite{liu2024vmamba}，本研究将其适配到EEG信号的生理特性，提出了ConvSS2D模块。如图~\ref{fig:C4-2}所示，该模块专门设计用于从EEG信号中进行分层时空特征提取，结合了浅层卷积网络\cite{borra2020interpretable}的初始特征提取能力与多方向S6块\cite{gu2023mamba}的高阶依赖建模能力。

\begin{figure}[htbp]
\centering
\includegraphics[width=0.7\linewidth]{C4-2ConvSS2D-CN.pdf}
\caption{ConvSS2D特征提取器设计}
\caption*{\centering (a) ConvSS2D块的前向传播流程；(b) SS2D工作流程。}
\label{fig:C4-2}
\end{figure}

ConvSS2D模块的输入是选择后的EEG信号$\mathbf{X} \in \mathbb{R}^{B \times C_l \times L}$，其中$B$表示批次大小，$C_l$表示EEG通道数，$L$表示时间维度。浅层卷积网络采用分层方式提取特征，首先通过核大小为$k_t$的时序卷积捕获局部时间模式，生成特征图：
\begin{equation}
\mathbf{X}_{\text{temp}}[b, f, c, t] = \sum_{k=1}^{k_t} \mathbf{X}[b, c, t+k-1] \cdot \mathbf{W}_{\text{temp}}[f, c, k],
\end{equation}
式中，$\mathbf{W}_{\text{temp}}[f, c, k]$表示时序卷积核，$b$表示批次索引，$f$表示滤波器索引，$c$表示通道索引，$t$表示时间索引。结果张量为$\mathbf{X}_{\text{temp}} \in \mathbb{R}^{B \times F \times C_l \times T_1}$，其中$F$表示时序滤波器数量，$T_1 = L - k_t + 1$表示卷积后的时间维度。
为了整合$C_l$个通道间的空间依赖关系，对$\mathbf{X}_{\text{temp}}$应用空间卷积，使用空间核$\mathbf{W}_{\text{spatial}}$：
\begin{equation}
\mathbf{X}_{\text{spatial}}[b, f, t] = \sum_{c=1}^{C_l} \mathbf{X}_{\text{temp}}[b, f, c, t] \cdot \mathbf{W}_{\text{spatial}}[f, c],
\end{equation}
式中，$\mathbf{W}_{\text{spatial}}[f, c]$是空间卷积核，结果$\mathbf{X}_{\text{spatial}} \in \mathbb{R}^{B \times F \times T_1}$聚合了空间信息，同时保持时间和滤波器维度。这种设计能够有效捕获EEG信号在空间分布上的相关性，特别是相邻电极间的协同活动模式。
空间卷积后的特征通过指数线性单元（Exponential Linear Unit，ELU）\cite{clevert2015fast}进行非线性激活，随后进行核大小为$k_p$的平均池化。池化操作定义为：
\begin{equation}
\mathbf{X}_{\text{pooled}}[b, f, t'] = \frac{1}{k_p} \sum_{t=1}^{k_p} \text{ELU}(\mathbf{X}_{\text{spatial}}[b, f, t+t'-1]),
\end{equation}
式中，$\mathbf{X}_{\text{pooled}} \in \mathbb{R}^{B \times F \times T_2}$表示时间压缩的特征图，$T_2 = T_1/k_p$是池化后的时间维度，$t'$是新的时间索引。ELU激活函数相比于传统的ReLU能够提供更平滑的梯度，有助于网络的稳定训练。
最后，特征通过$1 \times 1$卷积进行投影，将其转换到紧凑的嵌入空间。投影计算为：
\begin{equation}
\mathbf{X}_{\text{embed}}[b, t', e] = \sum_{f=1}^{F} \mathbf{X}_{\text{pooled}}[b, f, t'] \cdot \mathbf{W}_{\text{proj}}[e, f],
\end{equation} 
式中，$\mathbf{W}_{\text{proj}}[e, f]$是投影核，$e$是嵌入维度索引，结果$\mathbf{X}_{\text{embed}} \in \mathbb{R}^{B \times T_2 \times E}$表示在$E$维空间中每个时间段的嵌入表示。

为了充分利用时空依赖关系，建模这些时间嵌入段间的复杂相互关系是至关重要的。本章引入了二维选择性扫描（Selective Scan 2D，SS2D）模块\cite{liu2024vmamba}，通过沿多个遍历路径处理$\mathbf{X}_{\text{embed}}$进行分层建模。SS2D通过四个扫描方向顺序遍历二维嵌入：从左到右的行扫描（$\mathbf{S}_1$）、从上到下的列扫描（$\mathbf{S}_2$）、从左上到右下的对角扫描（$\mathbf{S}_3$）以及从右上到左下的对角扫描（$\mathbf{S}_4$）。
在每个方向$d$中，S6块对输入嵌入进行操作并输出处理后的序列$\mathbf{Y}_d$，计算为：
\begin{equation}
\mathbf{Y}_d = \text{S6}(\text{Scan}_d(\mathbf{X}_{\text{embed}}), \mathbf{A}_d, \mathbf{B}_d, \mathbf{C}_d, \Delta_d),
\end{equation}
式中，$\text{Scan}_d(\cdot)$根据扫描方向$d$重塑输入，$\mathbf{A}_d$是结构化状态矩阵，$\mathbf{B}_d$和$\mathbf{C}_d$是输入和输出投影矩阵，$\Delta_d$表示时变参数。S6块通过遵循递归关系更新特征，以线性计算复杂度高效建模沿每个扫描路径的长程依赖关系。处理完四个扫描路径的序列$\mathbf{Y}_d$后，SS2D模块使用具有可学习权重的交叉合并操作将这些输出融合到最终特征图$\mathbf{F}_\text{SS2D}$中：  
\begin{equation}
\mathbf{F}_\text{SS2D}(i,j) = \sum_{d=1}^4 \alpha_d \cdot \mathbf{Y}_d(i,j), 
\end{equation}
式中，$\sum_{d=1}^4 \alpha_d = 1$，$\mathbf{F}_\text{SS2D}(i,j)$表示位置$(i,j)$处的融合特征，$\alpha_d$是确定每个扫描路径重要性的可学习参数。这种多方向扫描和融合机制使得ConvSS2D能够捕获EEG信号中复杂的时空模式，为后续的分类和重构任务提供了丰富而判别性的特征表示。

\subsection{多组件复合损失函数}
为了实现通道选择与任务性能的协调优化，本章设计了统一的损失函数来联合训练CogSSM的任务无关和任务感知分支。该损失函数通过重构、分类和多样性损失的有机结合，确保准确的信号超分辨率、鲁棒的分类性能以及多样化的通道选择。该设计充分考虑了EEG信号分析的特殊需求，旨在信号保真度、任务性能和选择多样性之间建立平衡。

重构损失专为任务无关分支设计，用于从降维通道表示中恢复原始高密度EEG信号。该损失定义为原始信号$\mathbf{X}_{h}$与其重构对应物$\hat{\mathbf{X}}_{h}$之间的均方误差（Mean Squared Error，MSE）：
\begin{equation}
\mathcal{L}_{\text{recon}} = \frac{1}{B} \sum_{b=1}^{B} \frac{1}{C_h \cdot T} \|\mathbf{X}_{h} - \hat{\mathbf{X}}_{h}\|_2^2,
\end{equation}
式中，$B$表示批次大小，$C_h$表示高密度EEG通道数，$T$表示时间长度，$\mathbf{X}_{h}$表示原始高密度EEG信号，$\hat{\mathbf{X}}_{h}$表示重构的高密度EEG信号。重构损失确保了选择的通道能够保留足够的信息来重建原始信号的全局特征，这对于保持信号的神经生理学意义至关重要。通过最小化重构误差，网络学习到的通道选择策略能够捕获EEG信号中的主要成分和关键模式。

分类损失应用于任务感知分支，驱动模型区分任务特定的类别。该损失定义为真实标签$y_{b,k}$与预测类别概率$\hat{y}_{b,k}$之间的交叉熵损失：
\begin{equation}
\mathcal{L}_{\text{class}} = -\frac{1}{B} \sum_{b=1}^{B} \sum_{k=1}^{K} y_{b,k} \log \hat{y}_{b,k},
\end{equation}
式中，$B$表示批次大小，$K$表示类别总数，$y_{b,k} \in \{0,1\}$表示第$b$个样本第$k$类的真实标签，$\hat{y}_{b,k}$表示第$b$个样本第$k$类的预测概率。分类损失确保了选择的通道在特定任务上具有足够的判别能力，能够有效地区分不同的脑状态或认知任务。这种任务导向的优化使得网络能够学习到与特定应用场景最相关的通道组合。

为了避免通道选择的单一化倾向，本章引入多样性损失以促进批次间通道选择分数$\mathbf{s}_{:,c}$的变异性。该损失定义为所有通道选择分数标准差的负均值：
\begin{equation}
\mathcal{L}_{\text{diversity}} = -\frac{1}{C_h} \sum_{c=1}^{C_h} \text{std}(\mathbf{s}_{:,c}),
\end{equation}
式中，$C_h$表示高密度EEG通道数，$\mathbf{s}_{:,c}$表示第$c$个通道在整个批次中的选择分数向量，$\text{std}(\cdot)$表示标准差函数。多样性损失鼓励多样化的通道选择，避免冗余，提高表示容量。这一设计防止了网络过度依赖某些特定的通道组合，促进了更加鲁棒和泛化的通道选择策略。


总损失函数通过加权和整合上述三个组件，平衡重构保真度、任务特定准确性和通道选择多样性：
\begin{equation}
\mathcal{L}_{\text{sel}} = \mathcal{L}_{\text{recon}} + \lambda_{\text{class}} \cdot \mathcal{L}_{\text{class}} + \lambda_{\text{diversity}} \cdot \mathcal{L}_{\text{diversity}},
\end{equation}
式中，$\lambda_{\text{class}}$表示分类损失的权重系数，$\lambda_{\text{diversity}}$表示多样性损失的权重系数，这两个超参数用于控制分类和多样性损失对总损失的贡献程度。通过联合优化这一多目标损失函数，CogSSM有效地平衡了EEG信号保真度和任务特定性能，确保了对多样化应用的适应性，同时保持了鲁棒的双分支性能。

\section{实验设置与性能评估}
\subsection{数据集与基线方法设置}
本研究采用两个公开可获得的EEG图数据集OpenBMI MI\cite{lee2019eeg}和PhysioNet MI\cite{schalk2002general}对所提出的方法进行评估验证。OpenBMI MI和PhysioNet MI数据集均提供左手与右手运动想象任务的二元分类场景，具有标准化的实验协议和电极配置。两个数据集采用国际10-20电极系统的扩展版本，电极布局覆盖运动皮层、感觉皮层以及其他相关脑区，为运动想象信号的全面捕获提供了充足的空间覆盖。数据预处理流程包括8--30Hz带通滤波、工频干扰去除以及基线校正等标准步骤。

OpenBMI MI运动想象数据集记录了54名健康参与者的EEG信号，使用62通道Ag/AgCl电极系统，采样频率为1000Hz。Lee等人于2019年构建并公开发布了该数据集，目前已成为运动想象分类研究的标准基准之一。实验设计采用视觉提示范式，单个试次包含3秒准备期、4秒任务执行期（左手或右手运动想象）以及1.5-2.5秒休息期。受试者需要在两个独立会话中完成运动想象任务。本研究遵循已建立的评估协议\cite{lee2025neuroxai}，采用会话1数据进行模型训练（包括通道选择和分类器学习），采用会话2数据进行独立测试。这种时间分离的划分策略有效避免了数据泄露问题，确保了性能评估的客观性和可靠性。

PhysioNet运动想象数据集记录了109名受试者的64通道EEG信号，采样频率为160Hz。Schalk等人构建了该数据集，并通过PhysioBank平台进行公开发布\cite{goldberger_physiobank_2000}，已成为神经工程领域的标准数据资源。数据集涵盖运动执行和运动想象任务，本研究重点关注运动想象条件下的左手与右手二分类。实验数据以EDF+格式存储，包含标注完整的时间片段：休息状态（T0）、左手运动（T1）和右手运动（T2）。本研究采用严格的时间划分方案，奇数运行次数的数据用于训练，偶数运行次数的数据用于测试，保证了训练集和测试集之间的时间独立性。单个试次设计为4秒的运动想象周期，前置2秒的准备提示期，运行间隔设置为3-4分钟以减少疲劳效应。每位受试者完成14次运行，每次运行包含多个试次，为模型的训练和评估提供了充分的数据支撑。

为全面评估所提出方法的有效性，本研究构建了涵盖通道选择和运动想象分类的综合基线方法体系。表~\ref{tab:table1}汇总了各基线方法的核心特征，包括发表年份、开源可用性以及方法类型。根据实现原理，这些方法可分为特征工程（Feature Engineering, FE）和数据驱动（Data-Driven, DD）两大类别。


\begin{table}[htbp]
\centering
\caption{通道选择和分类的基线方法}
\begin{tabular}{ccccc}
\toprule
类别 & 方法 & 年份 & 开源 & 类型 \\ 
\midrule
\multirow{5}{*}{\begin{tabular}[c]{@{}c@{}}通道\\
选择器\end{tabular}} & Fisher \cite{tam2011minimal} & 2011 & 否 & FE \\
 & SVM-RFE \cite{feng2019multi} & 2019 & 是 & FE \\
 & CSP-RFE \cite{feng2019optimized} & 2019 & 是 & FE \\
 & Gumbel-softmax \cite{strypsteen2021end} & 2021 & 是 & DD \\
 & ECA \cite{tong2023learnable} & 2023 & 否 & DD \\ 
\midrule
\multirow{5}{*}{分类器} & CSP+LDA \cite{wang2008local} & 2008 & 是 & FE \\
 & FBCSP+LDA \cite{ang2008filter} & 2008 & 是 & FE \\
 & ShallowNet \cite{borra2020interpretable} & 2020 & 是 & DD \\
 & EEGNet \cite{lawhern2018eegnet} & 2018 & 是 & DD \\
 & Conformer \cite{song2022eeg} & 2021 & 是 & DD \\ 
\bottomrule
\end{tabular}
\begin{tablenotes}
\footnotesize
\item[]\centering 注：FE表示特征工程，DD表示数据驱动。
\end{tablenotes}
\label{tab:table1}
\end{table}

在通道选择方法方面，基线算法涵盖了特征工程和数据驱动的多种策略。Tam等人于2011年提出的Fisher方法\cite{tam2011minimal}基于统计分析原理，通过计算每个通道上特征的类间方差与类内方差比值来量化通道的判别能力，比值越高表明该通道对分类任务的贡献越大。Feng等人于2019年开发的CSP-RFE方法\cite{feng2019optimized}结合了共同空间模式和递归特征消除技术，首先利用全通道数据训练CSP模型，随后根据特征权重递归消除重要性最低的通道，直至达到目标通道数量。同期，Feng等人还提出了SVM-RFE方法\cite{feng2019multi}，该方法采用支持向量机执行类似的递归消除流程，通过分析SVM权重向量确定通道重要性的排序。Strypsteen等人于2021年引入的Gumbel-Softmax方法\cite{strypsteen2021end}采用端到端可学习框架，利用近似离散采样机制实现可微分的通道选择过程。Tong等人于2023年提出的ECA方法\cite{tong2023learnable}基于高效通道注意力机制，通过学习通道间的相互依赖关系，自适应调整各通道的重要性权重。

在分类算法方面，基线方法涵盖了从传统特征工程到现代深度学习的多种架构。Wang等人于2008年提出的CSP+LDA方法\cite{wang2008local}代表了经典特征工程范式，其核心在于寻找最优空间滤波器使得一类信号方差最大化而另一类方差最小化，然后采用线性判别分析完成分类。Ang等人于2008年开发的FBCSP+LDA方法\cite{ang2008filter}扩展了基础CSP框架，将EEG信号分解为多个频带，在各频带上独立应用CSP并融合所有频带特征，通过滤波器组策略处理多频带信息。现代数据驱动方法包括ShallowNet、EEGNet和Conformer等深度学习架构。Borra等人于2020年设计的ShallowNet\cite{borra2020interpretable}采用浅层卷积架构，专门针对EEG信号处理进行优化，在保持良好分类性能的同时具备较强的可解释性。Lawhern等人于2018年提出的EEGNet\cite{lawhern2018eegnet}是专为EEG图信号设计的紧凑卷积神经网络，通过深度卷积和分离卷积技术实现高效特征提取。Song等人于2021年开发的Conformer方法\cite{song2022eeg}融合了卷积神经网络和Transformer架构的优势，能够同时建模局部特征和全局依赖关系。

本研究提出的CogSSM方法通过自适应Gumbel-TopK机制实现通道选择功能，同时利用任务感知分支完成分类任务，构建了通道选择与运动想象分类的联合优化框架。该方法的核心创新在于实现了端到端的学习过程，能够自动发现最优通道子集和分类策略的协同配置。CogSSM的突出优势体现在其根据特定任务需求动态调整通道选择策略的能力。

\subsection{通道选择性能对比与分析}
\label{subsec:comparison_sota}
本研究在OpenBMI MI和PhysioNet MI两个数据集上全面评估了CogSSM方法的性能。采用平衡准确率和AUC-ROC指标进行全面验证，所有方法均在相同条件下运行，将原始通道集减少到20个电极，同时保持基于会话的受试者独立性。评估结果显示，CogSSM通过优化的通道选择策略和先进的ConvSS2D特征提取器实现了双重性能优势，在两个数据集上均表现出优于已建立基线方法的卓越性能。

\begin{table}[htbp]
\centering
\caption{不同方法通道选择后的平衡准确率对比}
\resizebox{\textwidth}{!}{%
\begin{tabular}{ccccccc}

\multicolumn{7}{c}{(a) OpenBMI MI数据集} \\
\toprule
方法 & CSP+LDA & FBCSP+LDA & ShallowNet & EEGNet & Conformer & 任务感知 \\ 
\midrule
Fisher & \uline{65.07 $\pm$ 10.90} & \uline{66.43 $\pm$ 13.42} & 69.80 $\pm$ 10.65 & 72.00 $\pm$ 12.75 & 70.61 $\pm$ 12.95 & 74.30 $\pm$ 12.89 \\
SVM-RFE & 63.94 $\pm$ 12.33 & 65.83 $\pm$ 10.93 & 68.61 $\pm$ 11.44 & 73.30 $\pm$ 13.40 & 71.67 $\pm$ 12.45 & \uline{74.81 $\pm$ 13.31} \\
CSP-RFE & 61.94 $\pm$ 9.31 & 63.44 $\pm$ 11.93 & 67.52 $\pm$ 9.96 & 68.11 $\pm$ 10.47 & 67.89 $\pm$ 8.06 & 72.48 $\pm$ 13.65 \\
Gumbel-Softmax & 63.44 $\pm$ 11.93 & 63.17 $\pm$ 9.31 & 70.65 $\pm$ 11.77 & \uline{74.00 $\pm$ 13.70} & \uline{73.09 $\pm$ 10.77} & 74.06 $\pm$ 13.88 \\
ECA & 62.00 $\pm$ 10.71 & 64.76 $\pm$ 11.74 & 69.09 $\pm$ 10.00 & 72.78 $\pm$ 10.45 & 70.50 $\pm$ 10.70 & 74.76 $\pm$ 14.06 \\
Manual & 65.44 $\pm$ 10.95 & 65.46 $\pm$ 12.33 & \uline{70.69 $\pm$ 10.85} & 72.76 $\pm$ 0.73 & 70.74 $\pm$ 13.36 & 73.98 $\pm$ 10.48 \\
CogSSM & \textbf{67.15 $\pm$ 13.32} & \textbf{71.65 $\pm$ 14.71} & \textbf{71.35 $\pm$ 12.42} & \textbf{74.26 $\pm$ 0.66} & \textbf{73.63 $\pm$ 13.06} & \textbf{76.95 $\pm$ 11.95} \\ 
\bottomrule
\multicolumn{7}{c}{(b) PhysioNet MI数据集} \\
\toprule
方法 & CSP+LDA & FBCSP+LDA & ShallowNet & EEGNet & Conformer & 任务感知 \\ 
\midrule
Fisher & \uline{68.98 $\pm$ 11.11} & 71.32 $\pm$ 9.10 & \uline{85.66 $\pm$ 7.44} & 80.79 $\pm$ 7.95 & 86.59 $\pm$ 5.70 & 86.04 $\pm$ 0.63 \\
SVM-RFE & 68.58 $\pm$ 8.48 & 70.66 $\pm$ 9.04 & 85.60 $\pm$ 7.23 & 79.69 $\pm$ 7.90 & \uline{86.80 $\pm$ 5.71} & 87.49 $\pm$ 5.53 \\
CSP-RFE & 68.44 $\pm$ 11.18 & 71.62 $\pm$ 10.38 & 83.78 $\pm$ 7.34 & 78.29 $\pm$ 7.89 & 86.40 $\pm$ 5.69 & 86.90 $\pm$ 5.77 \\
Gumbel-Softmax & 68.52 $\pm$ 11.57 & \uline{71.70 $\pm$ 9.63} & 84.96 $\pm$ 6.92 & 79.73 $\pm$ 7.96 & 86.32 $\pm$ 4.76 & 88.08 $\pm$ 5.09 \\
ECA & 66.02 $\pm$ 11.28 & 70.50 $\pm$ 9.89 & 85.12 $\pm$ 6.46 & \uline{80.90 $\pm$ 7.48} & 86.15 $\pm$ 5.90 & \uline{88.27 $\pm$ 5.51} \\
Manual & 66.75 $\pm$ 10.12 & 70.08 $\pm$ 8.57 & 84.26 $\pm$ 7.27 & 79.47 $\pm$ 8.03 & 85.91 $\pm$ 0.79 & 87.30 $\pm$ 5.84 \\
CogSSM & \textbf{69.34 $\pm$ 10.95} & \textbf{73.67 $\pm$ 8.81} & \textbf{88.51 $\pm$ 5.30} & \textbf{81.36 $\pm$ 8.06} & \textbf{87.75 $\pm$ 5.72} & \textbf{88.55 $\pm$ 7.30} \\ 
\bottomrule
\end{tabular}
}
\begin{tablenotes}
\footnotesize
\item[]\centering 注：手动选择的通道为FC-5/3/1/2/4/6，C-5/3/1/z/2/4/6，CP-5/3/1/z/2/4/6。最佳结果用\textbf{粗体}标示，次佳结果用\underline{下划线}标示。结果报告为平衡准确率的均值$\pm$标准差。
\end{tablenotes}
\label{tab:4-2}
\end{table}

\begin{table}[htbp]
\centering
\caption{不同方法通道选择后的AUC-ROC对比}
\resizebox{\textwidth}{!}{%
\begin{tabular}{ccccccc}

\multicolumn{7}{c}{(a) OpenBMI MI数据集} \\
\toprule
方法 & CSP+LDA & FBCSP+LDA & ShallowNet & EEGNet & Conformer & 任务感知 \\ 
\midrule
Fisher & \uline{0.7813 $\pm$ 0.1424} & \uline{0.7908 $\pm$ 0.1015} & 0.7271 $\pm$ 0.1264 & 0.7583 $\pm$ 0.1410 & 0.7341 $\pm$ 0.1564 & \uline{0.8046 $\pm$ 0.1419} \\
SVM-RFE & 0.7872 $\pm$ 0.1430 & 0.7863 $\pm$ 0.0964 & 0.7170 $\pm$ 0.1281 & 0.7743 $\pm$ 0.1462 & 0.7503 $\pm$ 0.1401 & 0.7725 $\pm$ 0.1370 \\
CSP-RFE & 0.7756 $\pm$ 0.1577 & 0.7735 $\pm$ 0.1356 & 0.7043 $\pm$ 0.1167 & 0.7319 $\pm$ 0.1178 & 0.7039 $\pm$ 0.1039 & 0.7560 $\pm$ 0.1163 \\
Gumbel-Softmax & 0.7702 $\pm$ 0.1590 & 0.7559 $\pm$ 0.1163 & 0.7381 $\pm$ 0.1375 & \uline{0.7911 $\pm$ 0.1430} & \uline{0.7704 $\pm$ 0.1152} & 0.7784 $\pm$ 0.1356 \\
ECA & 0.7809 $\pm$ 0.1587 & 0.7688 $\pm$ 0.1361 & 0.7212 $\pm$ 0.1177 & 0.7903 $\pm$ 0.1070 & 0.7418 $\pm$ 0.1209 & 0.7822 $\pm$ 0.1334 \\
Manual & 0.7729 $\pm$ 0.1151 & 0.7863 $\pm$ 0.0964 & \uline{0.7209 $\pm$ 0.1255} & 0.7893 $\pm$ 0.1464 & 0.8223 $\pm$ 0.1340 & 0.7820 $\pm$ 0.1348 \\
CogSSM & \textbf{0.7967 $\pm$ 0.0988} & \textbf{0.8183 $\pm$ 0.1467} & \textbf{0.7424 $\pm$ 0.1411} & \textbf{0.7941 $\pm$ 0.1197} & \textbf{0.7737 $\pm$ 0.1464} & \textbf{0.8897 $\pm$ 0.1066} \\ 
\bottomrule
\multicolumn{7}{c}{(b) PhysioNet MI数据集} \\
\toprule
方法 & CSP+LDA & FBCSP+LDA & ShallowNet & EEGNet & Conformer & 任务感知 \\ 
\midrule
Fisher & \uline{0.6861 $\pm$ 0.1583} & 0.7352 $\pm$ 0.1376 & \uline{0.8681 $\pm$ 0.0810} & 0.8347 $\pm$ 0.1069 & 0.8886 $\pm$ 0.0670 & 0.8650 $\pm$ 0.0849 \\
SVM-RFE & 0.6897 $\pm$ 0.1510 & 0.7333 $\pm$ 0.1362 & 0.8705 $\pm$ 0.0906 & 0.8269 $\pm$ 0.1014 & \uline{0.8824 $\pm$ 0.0817} & 0.9071 $\pm$ 0.0607 \\
CSP-RFE & 0.6779 $\pm$ 0.1485 & 0.7409 $\pm$ 0.1428 & 0.8570 $\pm$ 0.0977 & 0.8045 $\pm$ 0.0988 & 0.8839 $\pm$ 0.0749 & 0.9071 $\pm$ 0.0607 \\
Gumbel-Softmax & 0.6798 $\pm$ 0.1686 & \uline{0.7326 $\pm$ 0.1469} & 0.8499 $\pm$ 0.0896 & 0.8220 $\pm$ 0.1108 & 0.8768 $\pm$ 0.0693 & 0.9106 $\pm$ 0.0614 \\
ECA & 0.6546 $\pm$ 0.1611 & 0.7189 $\pm$ 0.1473 & 0.8588 $\pm$ 0.0866 & \uline{0.8332 $\pm$ 0.1077} & 0.8839 $\pm$ 0.0749 & \uline{0.9204 $\pm$ 0.0509} \\
Manual & 0.6570 $\pm$ 0.1457 & 0.7081 $\pm$ 0.1414 & 0.8570 $\pm$ 0.0977 & 0.8277 $\pm$ 0.1165 & 0.8863 $\pm$ 0.0769 & 0.9052 $\pm$ 0.0599 \\
CogSSM & \textbf{0.6954 $\pm$ 0.1639} & \textbf{0.7470 $\pm$ 0.1424} & \textbf{0.9030 $\pm$ 0.0850} & \textbf{0.8466 $\pm$ 0.0880} & \textbf{0.8888 $\pm$ 0.0702} & \textbf{0.9286 $\pm$ 0.0527} \\ 
\bottomrule
\end{tabular}%
}
\begin{tablenotes}
\footnotesize
\item[]\centering 注：手动选择的通道为FC-5/3/1/2/4/6，C-5/3/1/z/2/4/6，CP-5/3/1/z/2/4/6。最佳结果用\textbf{粗体}标示，次佳结果用\underline{下划线}标示。结果报告为AUC-ROC的均值$\pm$标准差。
\end{tablenotes}
\label{tab:4-3}
\end{table}


从通道选择策略的有效性角度分析，CogSSM展现出了卓越的通道识别和优化能力。在OpenBMI MI数据集上，如表~\ref{tab:4-2}(a)所示，当采用相同分类器架构时，CogSSM选择的通道始终产生更优的分类性能。在传统特征工程方法中，基于CogSSM选择通道的CSP+LDA分类器达到67.15\%的平衡准确率，相比次优通道选择方法Fisher的65.07\%提升了2.08\%。更为显著的是，在FBCSP+LDA分类器上，CogSSM选择的通道实现了71.65\%的准确率，相比最佳基线Fisher方法的66.43\%提升了5.22\%，这一显著改进充分证明了CogSSM通道选择策略在传统机器学习框架下的卓越效果。在深度学习分类器方面，CogSSM同样展现出了优异的通道选择能力。使用ShallowNet分类器时，CogSSM选择的通道达到71.35\%的准确率，相比最佳基线手动选择的70.69\%提升了0.66\%。对于EEGNet分类器，CogSSM实现了74.26\%的平衡准确率，不仅在均值上超越了最佳基线Gumbel-Softmax的74.00\%，更重要的是标准差显著降低，从13.70\%大幅减少至0.66\%，这表明CogSSM选择的通道具有更好的稳定性和跨受试者的一致性。在PhysioNet数据集上，如表~\ref{tab:4-2}(b)所示，CogSSM的通道选择优势更加明显。使用ShallowNet分类器时，CogSSM选择的通道达到88.51\%的准确率，相比最佳基线Fisher方法的85.66\%提升了2.85\%。这种显著的性能提升表明CogSSM能够识别出包含更丰富判别信息的通道子集，即使在相对简单的网络架构下也能实现优异的分类性能。对于传统方法CSP+LDA，CogSSM选择的通道实现了69.34\%的准确率，相比最佳基线Fisher的68.98\%提升了0.36\%。在FBCSP+LDA上的提升更为显著，达到73.67\%，相比最佳基线Gumbel-Softmax的71.70\%提升了1.97\%。

从ConvSS2D特征提取器的性能角度分析，CogSSM的任务感知分支充分展现了其在EEG信号模式识别方面的强大能力。在OpenBMI MI数据集上，CogSSM的任务感知分类器达到了76.95\%的平衡准确率，这一成绩超越了所有基线分类器的最佳表现。相比次优分类器SVM-RFE配合任务感知分支的74.81\%，CogSSM的任务感知分支提升了2.14\%。在AUC-ROC指标上，如表~\ref{tab:4-3}(a)所示，CogSSM任务感知分支的表现更加突出，达到0.8897，相比最佳基线Fisher配合任务感知分支的0.8046显著提升了8.51\%，这一提升充分证明了ConvSS2D特征提取器在EEG信号时空模式建模方面的卓越能力。在PhysioNet数据集上，CogSSM任务感知分支同样展现出了优异的性能表现，平衡准确率达到88.55\%，相比最佳基线ECA配合任务感知分支的88.27\%提升了0.28\%。在AUC-ROC指标上，如表~\ref{tab:4-3}(b)所示，CogSSM任务感知分支达到0.9286，相比最佳基线ECA配合任务感知分支的0.9204$\pm$0.0509提升了0.82\%。虽然PhysioNet数据集上的提升幅度相对较小，但考虑到该数据集上基线方法已经达到了相当高的性能水平，这种持续的性能改进仍然具有重要的实际价值。


CogSSM在不同分类器架构上展现的一致性性能提升验证了其设计的科学性和方法的普适性。无论是传统的特征工程方法如CSP+LDA、FBCSP+LDA，还是现代的深度学习方法如ShallowNet、EEGNet、Conformer，CogSSM选择的通道都能稳定实现性能提升。这种跨架构的一致性改进表明CogSSM能够识别出具有内在强判别性的通道组合，这些通道蕴含的神经生理信息能够被不同类型和复杂度的分类器有效捕获和利用。
综合以上全面分析可以得出，CogSSM方法的卓越性能源于两个核心技术的深度融合：第一，基于Gumbel-TopK的自适应通道选择机制通过可微分的离散优化过程，能够有效识别最具任务相关性和判别能力的电极子集，从源头上提升了输入信号的质量和信息密度；第二，ConvSS2D特征提取器通过其创新的多方向扫描机制和状态空间建模能力，能够深度挖掘EEG信号中的复杂时空依赖关系和非线性动态模式，为下游分类任务提供了更加丰富、抽象和判别性的特征表示。这种通道选择与特征提取的端到端协同优化策略使得CogSSM在EEG信号分析任务中展现出了卓越的性能表现和良好的跨数据集泛化能力。


\subsection{通道选择机制中的超参数优化}
\label{subsec:hyperparameter_optimization}

为深入理解通道选择模块的性能特征并获得最优配置，本研究采用网格搜索方法系统地探索关键超参数空间。超参数网格涵盖了影响通道选择性能的核心参数，包括历史权重系数（$\gamma$）、动量系数（$\alpha$）、探索率（$\varepsilon$）、温度参数（$\tau$）及其衰减率（$t$）。通过涵盖162种参数组合的系统性超参数优化实验，详细的参数配置如表~\ref{tab:C4-4}所示。


\begin{table}[h]
\centering
\caption{通道选择优化的超参数网格设置}
\label{tab:C4-4}
\begin{tabular}{ccc}
\toprule
参数 & 描述 & 取值 \\ \midrule
$\gamma$ & 历史权重系数 & \{0.0, 0.2, 0.4\} \\
$\alpha$ & 动量系数 & \{0.7, 0.8, 0.9\} \\
$\varepsilon$ & 探索率 & \{0.1, 0.3, 0.5\} \\
$\tau$ & 初始温度参数 & \{0.5, 1.0\} \\
$t$ & 温度衰减率 & \{0.03, 0.05, 0.07\} \\ \bottomrule
\end{tabular}
\end{table}


图~\ref{fig:C4-3}展示了不同超参数配置下的性能表现以及关键参数对模型性能的影响规律。通过对162种配置的全面评估，本研究识别出了性能最优的参数组合，并深入分析了各超参数的作用机制。


\begin{figure*}[htbp]
    \centering
    \includegraphics[width=0.9\linewidth]{C4-3hyper-CN.pdf}
    \caption{通道选择模块的超参数探索与性能分析}
    \caption*{\centering (a) 不同配置下的训练损失和验证损失变化趋势，其中配置ID 95表现出最优的验证性能；(b) 历史权重系数$\gamma$对验证损失的影响效果分析，展示了适度历史信息整合的重要性；(c) 动量系数$\alpha$对模型验证性能的影响，体现了长期稳定性机制的价值。}
    \label{fig:C4-3}
\end{figure*}


历史权重系数$\gamma$的实验结果揭示了该参数对模型性能的显著影响。如图~\ref{fig:C4-3}(b)所示，适度的历史信息整合（$\gamma=0.2$）能够有效降低验证损失至0.1362，相较于无历史信息条件（$\gamma=0.0$）的损失值0.1383，表现出明显的性能提升。历史权重系数的设计基于以下理论考量：通道选择过程应当借鉴先前选择的经验，以避免频繁的策略变动和不稳定的选择模式。当$\gamma=0.2$时，模型能够有效利用历史选择信息指导当前决策，这种适度的历史依赖性有助于平滑选择过程中的随机波动，提高选择策略的一致性和可靠性。
然而，进一步增加历史权重至$\gamma=0.4$导致验证损失上升至0.1405，这一现象表明过度依赖历史选择信息可能限制模型对输入数据变化的适应能力。当历史权重过高时，模型可能过分固化于先前的选择模式，缺乏对新信息和数据分布变化的敏感性。这种过度的历史依赖可能导致模型在面对不同受试者或实验条件时表现出较差的泛化性能。因此，历史权重系数的设置需要在稳定性和适应性之间寻求最优平衡，既要利用历史经验提高选择一致性，又要保持足够的灵活性应对数据变化。

动量系数$\alpha$的影响效果如图~\ref{fig:C4-3}(c)所示。实验结果表明，较高的动量值（$\alpha=0.9$）产生了最佳的性能表现，验证了通道选择过程中长期稳定性的重要价值。动量系数控制历史缓存$\mathbf{h}_t$的更新速率，较高的动量值使历史缓存能够在更长时间跨度内保留选择信息，有助于减轻批次间随机波动的影响，产生更加一致的选择模式。当$\alpha=0.9$时，历史缓存对过往信息的保留能力更强，能够有效抑制训练过程中的短期噪声和随机扰动。这种长期记忆机制使通道选择过程更加稳定，避免了因单个批次的异常数据或随机初始化导致的选择策略急剧变化。相比之下，较低的动量值（如$\alpha=0.7$或$\alpha=0.8$）虽然提供更快的适应能力，但可能导致选择过程过于敏感，容易受短期波动影响而产生不稳定的选择模式。

综合分析显示，最优配置（ID: 95）采用了适度的历史权重（$\gamma=0.2$）、高动量值（$\alpha=0.9$）、适度探索率（$\varepsilon=0.1$）以及合适的温度参数（$\tau=1.0$）和衰减率（$t=0.05$）。这种参数组合在选择稳定性和多样性之间实现了良好平衡，最终获得了0.10524的验证损失和0.14193的训练损失。
本节的超参数探索实验不仅验证了稳定化通道选择机制的有效性，也为参数优化提供了科学指导。通过历史信息整合和动量更新机制，成功地在选择稳定性和适应性之间建立了平衡，为EEG信号分析提供了可靠基础。

\subsection{消融实验与参数设置}
\label{subsec:ablation_analysis}

为了深入理解CogSSM框架各组件的作用机制并验证设计的合理性，本研究设计了两组系统性的消融实验。第一组实验重点分析了协同引导机制和状态空间建模的贡献度，第二组实验则探讨了多目标损失函数中各组件的作用机制。

（一）关键组件贡献度分析

为确保消融实验结果的可靠性和统计显著性，本研究从每个数据集中选择了20名表现出一致超随机性能的受试者，避免了BCI研究中常见的“运动想象文盲”现象对实验结果的干扰。所有实验均采用20个电极的固定配置，通过任务感知分支进行最终分类，确保了实验条件的一致性和结果的可比性。

\begin{table*}[htbp]
\centering
\caption{非BCI盲受试者的组件消融研究}
\label{tab:C4-5}
\resizebox{\textwidth}{!}{%
\begin{tabular}{ccccccc}
\toprule
\multirow{2}{*}{实验配置} & \multirow{2}{*}{\makecell{参数数量\\(M)}} & \multicolumn{2}{c}{OpenBMI MI} & \multicolumn{2}{c}{PhysioNet MI} \\ 
\cmidrule(lr){3-4} \cmidrule(lr){5-6}
 &  & 平衡准确率(\%) & AUC-ROC & 平衡准确率(\%) & AUC-ROC \\
\midrule
去除任务感知分支 & 2.07 & \uline{81.75 $\pm$ 3.85} & \uline{0.865 $\pm$ 0.028} & 87.88 $\pm$ 2.95 & 0.895 $\pm$ 0.022 \\
去除任务无关分支 & 2.05 & 79.50 $\pm$ 4.12 & 0.848 $\pm$ 0.041 & 86.36 $\pm$ 3.24 & 0.897 $\pm$ 0.031 \\
替换为Transformer & 2.38 & 80.00 $\pm$ 3.92 & 0.858 $\pm$ 0.037 & \uline{90.14 $\pm$ 2.68} & \uline{0.907 $\pm$ 0.019} \\
\midrule
CogSSM & 2.23 & \textbf{83.95 $\pm$ 2.76} & \textbf{0.904 $\pm$ 0.020} & \textbf{91.73 $\pm$ 2.15} & \textbf{0.927 $\pm$ 0.017} \\
\bottomrule
\end{tabular}%
}
\begin{tablenotes}
\footnotesize
\item[]\centering 注：最佳结果用\textbf{粗体}标示，次佳结果用\underline{下划线}标示。结果报告为均值$\pm$标准差。
\end{tablenotes}
\end{table*}

如表~\ref{tab:C4-5}所示，完整的CogSSM架构在两个数据集上始终优于所有消融变体。在OpenBMI MI数据集上，CogSSM达到了83.95\%的平衡准确率和0.904的AUC-ROC，相比最佳基线（仅任务无关分支）在准确率上提升了2.2\%，在AUC-ROC上提升了0.039。在PhysioNet MI数据集上，CogSSM达到了91.73\%的平衡准确率和0.927的AUC-ROC，相比最强基线（基于Transformer的架构）分别提升了1.59\%和0.020。

消融结果揭示了协同引导机制中两个分支的互补作用。当移除任务感知分支时，模型仅依赖重构损失进行优化，虽然能够保持信号的全局结构，但缺乏任务特异性的引导，导致选择的通道可能包含与分类任务无关的冗余信息。相反，当移除任务无关分支时，模型完全依赖分类损失，可能导致过度拟合到特定的分类任务，忽略了信号的全局结构和神经生理学完整性，这种单一目标的优化容易导致选择的通道在不同受试者或实验条件下缺乏泛化能力。此外，状态空间模型相对于Transformer架构的优势在实验中得到了充分验证。将基于SSM的特征提取替换为Transformer架构不仅使模型复杂度增加约6.7\%，同时产生较差的性能表现，这在OpenBMI数据集上尤为明显。这一发现证实了状态空间模型在捕获EEG信号固有复杂时序动态方面相对于注意力机制的效率和有效性。与Transformer中的全局注意力机制不同，SSM通过递归状态更新来建模序列依赖关系，这种机制更适合处理EEG信号的连续时序特性。在运动想象任务中，大脑的激活模式往往表现为复杂的时序展开过程，既包含从准备阶段到执行阶段的前向信息流，也包含从高级认知区域到运动皮层的反馈控制信号。ConvSS2D模块中实现的双向状态空间架构使模型能够更有效地建模神经信号中的因果和反因果关系。


（二）损失函数分析
\label{subsubsec:loss_analysis}

为了深入理解总损失函数$\mathcal{L}_{\text{sel}}$中各组件的作用机制，本研究对重构损失$\mathcal{L}_{\text{recon}}$、分类损失$\mathcal{L}_{\text{class}}$和多样性损失$\mathcal{L}_{\text{diversity}}$进行了系统的消融分析。如表~\ref{tab:C4-6}所示，通过比较不同损失组合下的性能表现，揭示了多目标优化策略的必要性。

\begin{table*}[htbp]
\centering
\caption{多目标损失函数的组件消融研究}
\label{tab:C4-6}
\begin{tabular}{ccccccc}
\toprule
\multicolumn{3}{c}{损失组件配置} & \multicolumn{2}{c}{OpenBMI MI} & \multicolumn{2}{c}{PhysioNet MI} \\ 
\cmidrule(lr){1-3} \cmidrule(lr){4-5} \cmidrule(lr){6-7}
$\mathcal{L}_{\text{recon}}$ & $\mathcal{L}_{\text{class}}$ & $\mathcal{L}_{\text{diversity}}$ & 平衡准确率(\%) & AUC-ROC & 平衡准确率(\%) & AUC-ROC \\
\midrule
$\checkmark$ &  &  & 77.25 $\pm$ 4.35 & 0.842 $\pm$ 0.025 & 85.61 $\pm$ 3.12 & 0.880 $\pm$ 0.020 \\
 & $\checkmark$ &  & 76.40 $\pm$ 4.82 & 0.835 $\pm$ 0.028 & 84.93 $\pm$ 3.45 & 0.875 $\pm$ 0.022 \\
$\checkmark$ & $\checkmark$ &  & \uline{83.17 $\pm$ 3.21} & \uline{0.892 $\pm$ 0.018} & \uline{91.71 $\pm$ 2.24} & \uline{0.923 $\pm$ 0.015} \\
$\checkmark$ &  & $\checkmark$ & 79.53 $\pm$ 4.12 & 0.852 $\pm$ 0.023 & 87.28 $\pm$ 2.98 & 0.885 $\pm$ 0.019 \\
 & $\checkmark$ & $\checkmark$ & 78.91 $\pm$ 4.26 & 0.848 $\pm$ 0.024 & 86.75 $\pm$ 3.15 & 0.882 $\pm$ 0.021 \\
\midrule
$\checkmark$ & $\checkmark$ & $\checkmark$ & \textbf{83.45 $\pm$ 3.05} & \textbf{0.894 $\pm$ 0.017} & \textbf{91.86 $\pm$ 2.18} & \textbf{0.925 $\pm$ 0.014} \\
\bottomrule
\end{tabular}%
\begin{tablenotes}
\footnotesize
\item[]\centering 注：最佳结果用\textbf{粗体}标示，次佳结果用\underline{下划线}标示。结果报告为均值$\pm$标准差。
\end{tablenotes}
\end{table*}

消融实验结果揭示了不同损失函数组件在优化通道选择性能中的协同作用机制。从表~\ref{tab:C4-6}可以观察到，各个损失组件均对最终的分类性能产生积极影响，且不同组件间存在显著的协同效应。仅使用重构损失的配置在OpenBMI MI数据集上获得了77.25\%的平衡准确率和0.842的AUC-ROC值，而仅使用分类损失的配置获得了相对较低的76.40\%平衡准确率和0.835的AUC-ROC值。在PhysioNet MI数据集上，重构损失配置达到了85.61\%的平衡准确率和0.880的AUC-ROC值，分类损失配置获得了84.93\%的平衡准确率和0.875的AUC-ROC值。这一结果表明，单纯依赖信号重构或任务分类的单一目标优化均无法充分发挥通道选择的潜力。
重构损失与分类损失的联合优化显著提升了模型性能。在OpenBMI MI数据集上，双重损失配置相较于单一重构损失提升了5.92\%的平衡准确率（从77.25\%提升至83.17\%），相较于单一分类损失提升了6.77\%。在PhysioNet MI数据集上，双重损失配置相较于单一重构损失和分类损失分别提升了6.10和6.78\%的平衡准确率。同时，AUC-ROC指标也呈现出相似的提升趋势，进一步验证了双重约束机制的有效性。

重构损失的核心作用在于维护信号的神经生理学完整性和全局结构特征。通过最小化原始高密度EEG信号与通过选择通道重构信号之间的均方误差，重构损失确保了选择的通道子集能够最大程度地保留原始信号的关键信息和空间分布特性。这种全局约束机制不仅防止了通道选择过程中的信息损失，更重要的是，它作为一种有效的正则化手段，避免了模型在特定分类任务上的过度拟合，从而提高了模型的泛化能力。分类损失则专注于优化任务相关的判别特征提取能力。通过交叉熵损失函数的引导，分类损失使得通道选择过程能够识别并优先选择那些对目标认知任务具有高判别性的电极位置。这种任务导向的优化策略确保了选择的通道在特定的BCI任务中能够有效区分不同的脑状态模式，从而实现优异的分类性能。分类损失的引入使得通道选择不再仅仅考虑信号的完整性保持，而是同时兼顾了任务特异性的判别需求。双重约束优化策略的成功验证了本研究的核心假设：有效的EEG通道选择需要在信号保真度和任务判别性之间实现平衡。这种平衡机制使得选择的通道既能保持原始信号的神经生理学意义，又能在特定认知任务上展现出卓越的分类能力。

多样性损失的加入进一步优化了模型性能和稳定性。在OpenBMI MI数据集上，完整的三重损失配置相较于双重损失配置额外提升了0.28\%的平衡准确率，在PhysioNet MI数据集上实现了0.15\%的提升。更为重要的是，多样性损失的引入显著改善了模型的稳定性，这体现在跨受试者性能变异性的明显降低上。在OpenBMI MI数据集上，平衡准确率的标准差从3.21\%降低至3.05\%，在PhysioNet MI数据集上从2.24\%降低至2.18\%。
多样性损失的正则化机制基于EEG信号空间分布的互补性原理。不同通道记录的神经活动在空间上具有互补特性，选择空间分布合理且信息互补的通道组合能够最大化整体信息获取量。多样性损失通过在批次训练过程中鼓励通道选择分数的变异性，促使网络探索多样化的通道组合策略，避免过度依赖特定的电极配置模式。这种探索机制不仅提高了通道选择的信息丰富度，还增强了模型对个体差异和环境噪声的鲁棒性。

实验结果充分验证了本研究提出的多目标优化框架的理论合理性和实践有效性。有效的EEG通道选择需要同时优化三个关键目标：信号保真度、任务判别性和表示多样性。重构损失通过维护全局信号结构防止了任务特异性过拟合；分类损失确保了对目标认知范式的高效判别能力；多样性损失作为正则化机制促进了信息丰富且非冗余的电极选择，提高了模型对噪声扰动和受试者间变异的适应能力。


\section{结果讨论}
\label{sec:discussion}

\subsection{通道选择机制的神经关联性分析}
如图~\ref{fig:C4-4}所示，本研究通过对训练过程中通道选择频率分布的深入分析，揭示了CogSSM模型所学习的通道选择策略远超传统运动想象电极配置的复杂性和分布特征，为理解运动想象任务的神经基础提供了新的视角。

\begin{figure}[!ht]
    \centering
    \includegraphics[width=0.94\linewidth]{C4-4频率-CN.pdf}
    \caption{运动想象通道选择的时空分布特征}
    \label{fig:C4-4}
    \caption*{\centering (a) 各通道在整个训练过程中的选择频率分布，红色虚线表示50\%选择阈值，蓝色点线表示平均选择频率；(b) 训练过程中通道选择模式的时间演化，颜色深度表示通道被选择的状态。}
\end{figure}

从图~\ref{fig:C4-4}(a)的选择频率分布可以观察到，CogSSM模型识别出的关键通道呈现显著的分布式特征。选择频率最高的通道（选择频率>0.9）主要包括FC1、C1、CP5、CP3、C2、C4等，这些通道构成了一个跨越多个功能脑区的神经网络。值得注意的是，虽然传统的感觉运动区域（如C3、C4、Cz）确实被高频选择，但模型同时一致性地纳入了来自额叶（F3、F1、F6、Fpz）、颞叶（FT8）、顶叶（P7、P3、Pz、P2）和枕叶区域（PO7、PO3、Oz、O2、Iz）的通道。这种选择模式表明，有效的运动想象解码需要整合来自分布式皮层网络的信息，而非仅依赖传统定义的感觉运动中心区域。图~\ref{fig:C4-4}(b)展示的通道选择时间演化揭示了模型学习过程的渐进性和稳定性特征。在训练初期（epochs 0-10），通道选择呈现相对随机的探索模式；随着训练的深入（epochs 10-30），核心感觉运动区域的选择逐渐稳定；在训练后期（epochs 30-60），完整的分布式网络模式逐步形成并趋于稳定。这种渐进式的网络构建过程表明，CogSSM模型能够通过自适应学习机制逐步完善对运动想象任务相关神经动力学的理解，从而实现从粗粒度到细粒度的特征提取优化。

分布式选择模式的神经生理学基础可以从多个层面进行解释。首先，额叶电极的高频选择反映了执行控制和运动规划网络在运动想象任务中的关键作用。前额叶皮层作为执行功能的核心区域，在复杂认知任务的注意分配、工作记忆维持和目标导向行为调控中发挥重要作用\cite{hanakawa2008motor}。其次，顶叶区域的一致性选择揭示了空间注意和身体图式在运动想象过程中的重要贡献。顶叶皮层不仅参与感觉运动整合，还负责身体空间表征和运动意图的编码，这对于产生清晰、稳定的运动想象至关重要。枕叶视觉区域电极的频繁选择可能反映了运动想象任务中视觉-空间处理的重要性。现代认知神经科学研究表明，运动想象过程往往伴随着丰富的视觉表象，包括对肢体运动轨迹、空间位置和动作场景的内在视觉化\cite{henschke2023engaging}。枕叶区域的激活可能捕获了这种内在视觉表象的神经基础，或者反映了与实验范式中视觉线索处理相关的神经活动。这一发现对运动想象BCI的设计具有重要的方法学启示：标准的运动想象协议通常采用视觉方向线索来引导受试者的想象内容，这可能会在解码过程中引入与视觉诱发电位相关的分类信息。

与预先定义感觉运动区域进行运动想象解码的传统方法相比，本研究提出的数据驱动通道选择机制展现出更强的适应性和个体化能力。这种方法不仅能够自动识别任务相关的关键通道，还能够适应不同个体间的神经解剖差异和功能网络变异。

\subsection{通道选择对不同分类架构的差异化影响}
基于表~\ref{tab:4-2}和表~\ref{tab:4-3}的实验结果，本研究发现通道选择策略对基于特征工程和基于深度学习的分类方法产生了显著的差异化影响，这种差异反映了两类方法在信息处理机制上的本质区别。

传统的基于特征工程方法在通道选择后表现出显著的性能改善，这一现象可从其固有的信息处理局限性进行解释。特征工程方法依赖于预定义的特征提取算子（如功率谱密度、共空间模式等），这些算子对输入信号的质量和相关性具有高度敏感性。当包含大量冗余或噪声通道时，不相关信息会在特征空间中引入混淆因子，降低类间可分离性。通过精确的通道选择，冗余信息得到有效过滤，使得人工设计的特征算子能够在更纯净的信号空间中提取判别性信息，从而显著提升分类性能。
相比之下，基于深度学习的方法对通道选择表现出较低的依赖性，这源于其分层表示学习和注意力机制的内在优势。深度神经网络通过多层非线性变换能够自动学习输入特征的重要性权重，有效实现隐式的通道选择和特征融合。这种自适应能力使得深度学习方法能够在包含噪声通道的情况下仍保持相对稳定的性能，体现了端到端学习范式在处理复杂、高维数据时的鲁棒性优势。

为了深入量化这种差异化影响，本研究设计了渐进式通道扩展实验：从性能最优的5个通道开始，逐步以5个通道为单位随机添加非最优通道，系统性地观察两类方法的性能退化模式。图~\ref{fig:C4-5}的实验结果揭示了两类方法在面对通道数量变化时的本质性差异：基于特征工程方法的性能曲线呈现明显的倒U型分布，在最优通道数量（约15-25个通道）达到峰值后迅速下降。这种陡峭的性能退化反映了传统方法对精心筛选通道的强烈依赖性——额外的噪声通道会直接干扰预定义特征的提取效果，导致特征空间的维数灾难和过拟合风险。相比之下，基于深度学习方法的性能曲线表现出更强的平稳性，即使在通道数量显著增加时仍能维持相对稳定的性能水平。深度学习方法的鲁棒性源于其多层次的特征学习架构。在网络的浅层，卷积和池化操作能够提取局部时频特征；在中间层，注意力机制和残差连接实现跨通道的信息融合和权重调节；在深层，全连接层完成高级语义特征的抽象。这种分层处理机制使得网络能够自动识别和强化有用信号，同时抑制噪声通道的干扰\cite{altaheri2023deep, gao2022convolutional}。

\begin{figure}[!ht]
    \centering
    \includegraphics[width=\linewidth]{C4-5methods-CN.pdf}
    \caption{不同分类架构在渐进式通道扩展下的性能变化对比}
    \label{fig:C4-5}
    \caption*{\centering (a) OpenBMI数据集上的性能变化曲线；(b) PhysioNet数据集上的性能变化曲线。蓝色箭头标识性能下降的临界点。}
\end{figure}

这种差异化特征为不同应用场景下的技术选型提供了重要指导。
（1）资源受限场景：在便携式或可穿戴BCI系统中，硬件成本和功耗是关键约束。此时，精确的通道选择与轻量级特征工程方法的结合能够实现性能与效率的最优平衡，通过最小化电极数量降低系统复杂度的同时保持竞争性的分类精度；
（2）高精度应用场景：在临床诊断或高精度研究平台中，当计算资源充足且最大分类准确率是首要目标时，基于深度学习的方法展现出更强的适用性。其隐式通道选择和自适应特征学习能力能够充分利用多通道信息，实现更高的分类性能上限；
（3）数据稀缺场景：在训练样本有限或信噪比较低的条件下，即使深度学习方法也能从显式的通道选择中获益。通过预先去除明显的噪声通道，可以减少模型的学习负担，提高训练效率和泛化能力。
这种基于方法学特性的适应性选择策略为BCI系统的工程化部署提供了理论指导，有助于在不同应用需求和资源约束下实现最优的系统设计。

\subsection{通道数量与质量之间的权衡分析}
在BCI系统的实际应用中，通道数量的选择往往面临性能提升与实用性约束之间的复杂权衡。一方面，更多的通道数量理论上能够提供更丰富的神经信息，有利于提高分类精度；另一方面，过多的通道会增加系统复杂性、计算负担和用户使用成本。

图~\ref{fig:C4-6}展示的验证损失与电极数量关系曲线揭示了通道选择中的边际效用递减规律。该曲线呈现明显的非线性特征，可划分为三个典型阶段：首先，在电极数量从1增加至11的过程中，验证损失出现急剧下降，表明核心通道的引入能够显著提升模型的分类能力；随后，当电极数量超过11时，损失下降趋势明显放缓，呈现渐进式改善模式；最终，在电极数量达到约20个后，损失函数趋于平稳甚至略有反弹，显示出通道冗余和噪声干扰的负面影响。
这种三阶段变化模式深刻反映了神经信息处理的内在机制。在初始阶段，每增加一个关键通道都能引入新的判别性神经特征，显著增强模型对不同运动想象任务的区分能力。然而，随着通道数量的进一步增加，新增通道所携带的独特信息逐渐减少，更多地表现为已有信息的重复或变异。当通道数量超过某个临界阈值后，边际信息增益急剧下降，而噪声累积和计算复杂度却持续增加，导致性能改善的成本效益比显著恶化。

\begin{figure}[ht]
    \centering
    \includegraphics[width=0.85\linewidth]{C4-6balance-CN.pdf}
    \caption{CogSSM模型验证损失随电极数量变化的优化曲线}
    \label{fig:C4-6}
\end{figure}

基于$K=11$这一平衡点的深入分析，本研究发现该配置在多个维度上实现了最优权衡。从性能角度，11个精选通道能够捕获运动想象分类所需的核心神经特征，验证损失降至1.994的相对低值；从实用性角度，该通道数量在保证分类精度的同时，显著降低了电极安装的复杂性和时间成本；从计算效率角度，适中的通道数量有效控制了数据维度，减少了实时处理的计算负担。这一发现为不同应用场景下的系统设计提供了量化依据。在便携式消费级BCI产品中，精简的电极配置直接转化为成本优势和用户体验的改善。

\section{本章小结}
本章提出的CogSSM协同引导学习框架通过将认知神经科学中的双重处理机制转化为可计算的算法架构，有效解决了便携式EEG系统中通道选择优化的核心难题。该框架采用全局整合和任务特化的协同优化策略：任务无关路径通过超分辨率重建确保所选通道子集能够保持原始高密度信号的统计一致性和频谱完整性，实现对分布式神经信息的全局保护；任务相关路径则通过分类目标优化提升特定应用的判别性能，实现对任务关键特征的精准提取。在技术创新方面，ConvSS2D模块融合了二维卷积的局部特征捕获能力与状态空间模型的长程依赖建模优势，在二维状态空间内联合刻画EEG信号的时空动力学模式。实验验证表明，该方法在OpenBMI与PhysioNet MI数据集上仅使用20个通道就分别达到了76.95\%与88.55\%的分类准确率，相比竞争基线在平衡准确率上提升2.14\%-15.01\%，AUC-ROC指标最高提升8.5\%，有效实现了系统复杂度的显著降低与任务性能的同步提升。
通道选择结果的深入分析揭示了超越传统神经解剖学区域划分的重要发现。CogSSM所选择的电极配置并非严格遵循基于感觉运动皮层的先验假设，而是体现了涵盖额叶、顶叶、颞叶和枕叶的复杂分布式皮层网络参与模式。这一发现不仅为运动想象等BCI任务的神经机制理解提供了新的见解，更为便携式EEG系统的电极布局设计提供了数据驱动的科学依据。渐进式通道扩展实验进一步揭示了不同分类架构对通道选择的差异化响应机制：基于特征工程的方法对通道质量表现出高度敏感性，而基于深度学习的方法则展现出更强的鲁棒性和自适应能力。

从技术体系构建的角度来看，CogSSM与第三章提出的MASER方法形成了增强--约简的协同技术方案。MASER从空间分辨率增强维度解决低通道数据的高保真重建问题，CogSSM从通道配置约简维度面向部署约束提供优化的通道布局方案。因此，在信号获取阶段通过CogSSM实现智能化的通道选择，在信号处理阶段通过MASER实现高质量的空间重建，共同构建了多层级生成建模框架中不可或缺的底层技术支撑。
